\begin{abstract}
Métodos de \textit{steering}, ou direcionamento das ativações durante a inferência, prometem reduzir a menção a objetos ausentes em modelos de visão e linguagem, mas seus resultados dependem do protocolo e podem refletir mudanças no comportamento de resposta. Avaliamos VTI e REVIS sem ajuste de hiperparâmetros em LLaVA-1.5-7B, modelo contemplado pelas publicações, e em LLaVA-OneVision-1.5-8B e Qwen3-VL-8B, não contemplados por elas. Comparamos receitas oficiais e controles aleatórios em POPE adversarial, RePOPE, CHAIR e MME, considerando propensão a responder sim, discriminação, cobertura de objetos, comprimento e repetição. Não se estabelece melhora robusta em POPE além dos controles. As reduções de CHAIR são limitadas por perda de cobertura ou degeneração das descrições; em OneVision, REVIS reduz CHAIR por descrição de 33,4\% para 25,6\%, mas apenas 1,9 ponto percentual excede o atalho de truncamento com \textit{recall} (cobertura de objetos) equivalente. Diagnósticos convergem para uma rota da imagem à palavra do objeto na pergunta e, depois, à posição que prediz a resposta, mais tardia nos modelos baseados em Qwen. Os achados mostram por que intervir na representação não implica melhorar a discriminação e por que avaliar steering exige controles de comportamento e de capacidade.
\end{abstract}
\begin{IEEEkeywords}
alucinação de objetos, modelos de visão e linguagem, steering, interpretabilidade, avaliação.
\end{IEEEkeywords}
